\section{Progressive Calibration Ladder}
\label{sec:calibration}

\subsection{Why a Ladder Rather Than a Switch}

Comparing an uncalibrated simulator against a calibrated one answers a
question nobody asked: it reports that fidelity changed something, without
saying which mismatch mattered. Since the goal is to make the simulator
usable as an evaluation instrument, the useful question is which \emph{class}
of mismatch controls its predictivity, because that identifies where
modelling effort should go. We therefore define four levels in which each
rung adds one class of measured calibration and explicitly leaves the
remaining classes untouched (Table~\ref{tab:calibration}). Declaring what is
left uncalibrated is as important as declaring what is calibrated: it is what
licenses attributing a change in predictivity to the rung that changed.

\subsection{Levels}

\textbf{\SZERO{} --- baseline simulation.} The default practical model: the
simulator's own camera configuration, no injected latency, and vehicle
controller gains and limits chosen for stable simulated behaviour rather than
measured against hardware. The step responses in
Table~\ref{tab:stepresponse} characterize this level. \SZERO{} is not a straw
man; it is what a competent practitioner obtains before any measurement
campaign, and it is the level most published underwater simulation results
implicitly occupy.

\textbf{\SONE{} --- camera and perception.} The simulated camera is matched
to the measured physical camera: resolution, frame rate, field of view and
intrinsics, and such image statistics as can be measured and reproduced.
Timing, vehicle response, and sensor noise remain at \SZERO{}. Following
Elmquist and Negrut~\cite{elmquist2022aperform}, the camera model is
validated by the response of the downstream detector rather than at pixel
level, because the detector's behavior is what the avoidance stack actually
consumes. \CALIB

\textbf{\STWO{} --- camera plus system timing.} Everything from \SONE{},
plus the measured end-to-end delay budget: acquisition, encoding and
transport, detector latency, and command-path latency, together with the
measured frame-rate and frame-drop behavior. The simulator injects these
through the command-latency buffer and the image path rather than through a
lumped constant, so the individual contributions remain separable. Vehicle
response and sensor noise remain at \SZERO{}. This level is the reason the
simulation clock had to be made honest (Sec.~\ref{sec:system}): injecting a
50\,ms delay is meaningless if simulated time and physics time already
disagree. \CALIB

\textbf{\STHREE{} --- camera, timing, and vehicle response.} Everything from
\STWO{}, plus the measured surge, sway, and yaw responses of the physical
vehicle --- time constants, steady-state gains, deadbands, and rate limits ---
and measured noise and update-rate models for attitude and depth. The
identification uses the same step-response script that produced
Table~\ref{tab:stepresponse}, so the two domains are characterized by an
identical procedure. Hydrodynamic identification of BlueROV2-class vehicles
is itself well
established~\cite{benzon2022anopenso,orjales2025towardsa}; the contribution
here is not the identification but its placement as one isolated rung of a
predictivity experiment. \CALIB

%==============================================================================
% S0-S3 calibration ladder.  Status: DRAFT (methodology fixed; the measured
% values that populate each level do not exist yet -- see paper/STATUS.md).
%==============================================================================
\begin{table}[t]
\centering
\caption{Calibration ladder. Each level adds one class of \emph{measured}
calibration; the ``left uncalibrated'' column is what licenses attributing a
change in predictivity to the rung that changed. No level is tuned against
avoidance outcomes.}
\label{tab:calibration}
\footnotesize
\setlength{\tabcolsep}{3pt}
\renewcommand{\arraystretch}{1.2}
\begin{tabular}{@{}l p{27mm} p{27mm}@{}}
\toprule
\textbf{Level} & \textbf{Calibrated against measurement} &
\textbf{Deliberately left uncalibrated} \\
\midrule
\SZERO   & --- (default practical model) & camera, timing, vehicle response, sensor noise \\
\SONE    & camera resolution, frame rate, FOV/intrinsics, measurable image statistics & timing, vehicle response, sensor noise \\
\STWO    & \SONE{} $+$ acquisition/encode/transport latency, detector latency, command latency, frame drops & vehicle response, sensor noise \\
\STHREE  & \STWO{} $+$ surge/sway/yaw response, deadbands, rate limits, attitude/depth noise and rate & optical medium (scattering, turbidity), tether dynamics \\
\bottomrule
\end{tabular}
\end{table}


\subsection{What the Ladder Deliberately Does Not Do}

Three exclusions keep the interpretation clean. First, no level tunes the
simulator against \emph{avoidance outcomes}: every calibrated quantity is
measured independently against its physical counterpart, which distinguishes
this design from post-hoc parameter search against the target
metric~\cite{kadian2020simrealp} and prevents the ladder from fitting the
answer. Second, no level changes the autonomy code, so a change in outcome
cannot come from a re-tuned controller. Third, the underwater optical medium
is calibrated only to the extent that it is measured at \SONE{}; scattering
and turbidity models are outside the scope of this study and remain a
declared limitation.

We also pre-register the interpretation of a negative result. Truong et
al.~\cite{truong2022rethinki} showed that added fidelity can degrade
sim-to-real behavior, so ``calibration at level $S_k$ did not improve
predictivity'' is a reportable finding of this experiment and not a failure of
it. The prediction being tested is directional, not guaranteed: we expect
timing and vehicle response to matter more than nominal camera geometry for
a commitment-based avoidance policy, and the experiment is designed to be
able to contradict that expectation.
